% BANKS-SOURCE: pdf=99 print=89 kind=problem-section id=chapter07-problems
\begin{banksproblems}

% BANKS-SOURCE: pdf=99 print=89 kind=problem id=7.1
\BanksProblem{7.1}{*}
Use Noether's theorem to write the conserved currents associated with the
$O(n)$ or $U(n)$ symmetry of a general invariant Lagrangian with up to two
derivatives, for scalar fields transforming in the fundamental (defining)
representation of these groups. Show that the Lagrangian for massless Dirac
fermions transforming in $N_f$ copies of a representation $R$ of any Lie group
$G$ has a $U(N_f)\times U(N_f)$ symmetry of separate unitary transformations on
left- and right-handed components of the Dirac fermions. We will see that most
of this symmetry remains valid when we couple the fermions to vector fields via
\begin{equation*}
  A_\mu^a\,\bar\psi\gamma^\mu t^a\psi,
\end{equation*}
The chiral and anomaly statements here are four-dimensional, $D=4$, with
$\{\gamma^\mu,\gamma^\nu\}=-2\eta^{\mu\nu}$ and
$\gamma_5=\ii\gamma^0\gamma^1\gamma^2\gamma^3$, and
$\slashed p=\gamma^\mu p_\mu$.  Here $t^a$ are the representatives of
infinitesimal $G$ transformations in
the $R$ representation. A certain $U(1)$ subgroup of
$U(N_f)\times U(N_f)$ is broken by a subtle quantum effect called an anomaly.

% BANKS-SOURCE: pdf=100 print=90 kind=problem id=7.2
\BanksProblem{7.2}{}
Use Noether's theorem to evaluate the conserved current corresponding to
translation symmetry, for general scalar field theories coupled to the Maxwell
field (couple them via the minimal substitution principle:
$\partial_\mu\phi_i\to(\partial_\mu-\ii e q_i A_\mu)\phi_i$; $e$ is the
coupling constant and $q_i$ the charge on the $i$th complex field). You should
find the momentum is the integral of the time component of the stress-energy
tensor:
\begin{equation*}
  P^\nu=\int\dd^{D-1}\mathbf x\,T^{0\nu},
  \qquad \partial_\mu T^{\mu\nu}=0.
\end{equation*}
The $T_{\mu\nu}$ you find will not be symmetric. Another way to define
$T_{\mu\nu}$ is to construct the action in a varying space-time metric (for
the Lagrangian we are describing this is simple: simply replace the Minkowski
metric $\eta_{\mu\nu}$ by $g_{\mu\nu}(x)$ everywhere, and replace
$\dd^D x\to\dd^D x\sqrt{-g}$, where $g$ is the determinant of the metric). The
stress tensor is defined by
\begin{equation*}
  T_{\mu\nu}(x)=-\frac{2}{\sqrt{-g}}\frac{\delta S}{\delta g^{\mu\nu}(x)},
\end{equation*}
and is obviously symmetric (we take $g_{\mu\nu}\to\eta_{\mu\nu}$ after taking
the derivative). Show that the two definitions give the same value for
$P_\mu$ as long as fields fall off sufficiently rapidly at infinity. Using the
gravitational definition define
\begin{equation*}
  M_{\mu\nu\lambda}=x_\mu T_{\nu\lambda}-x_\nu T_{\mu\lambda}.
\end{equation*}
Show that $\partial^\lambda M_{\mu\nu\lambda}=0$, and that
$J_{\mu\nu}=\int\dd^{D-1}\mathbf x\,M_{\mu\nu0}$ is the conserved angular-momentum
generator.

% BANKS-SOURCE: pdf=100 print=90 kind=problem id=7.3
\BanksProblem{7.3}{*}
Work in $D=4$ for this conformal-symmetry problem.
Suppose that the stress tensor is traceless:
\begin{equation*}
  \eta^{\mu\nu}T_{\mu\nu}=0.
\end{equation*}
Show that there is a five-parameter set of quadratic polynomials
\begin{equation*}
  f^\mu=Cx^\mu+C^\mu{}_{\nu\kappa}x^\nu x^\kappa
\end{equation*}
(you must find the allowed form of $C^\mu{}_{\nu\kappa}$) such that
$f^\mu T_{\mu\nu}$ is conserved. Together with the Poincare generators, the
corresponding Noether charges form the Lie algebra of the group of conformal
transformations $SO(2,4)$. Field theories invariant under these transformations
are called conformal field theories (CFTs). In our discussion of renormalization
we will see that the extreme ultraviolet and infrared behavior of any field
theory is described by a pair of conformal field theories (though the IR limit
may be trivial, and have only delta-function correlations). Any field theory
may be constructed by perturbation theory around its UV conformal limit.

% BANKS-SOURCE: pdf=101 print=91 kind=problem id=7.4
\BanksProblem{7.4}{*}
The pion field $\pi^a(x)$ transforms as a triplet (vector) of the $SU(2)$
isospin symmetry of strong interactions. Show that, in the limit in which we
consider this symmetry to be exact, the Ward identities for Green functions,
and the LSZ formula, imply that all three pions have the same mass. Now use
the symmetry and the obvious invariance of the four-point function
$\langle\pi^{a_1}(x_1)\cdots\pi^{a_4}(x_4)\rangle$ under permutation of the pion
fields to show that the $3^4$ different components of this Green function are
related to a much smaller number of functions of $x_1\to x_4$. Use the LSZ
formula to conclude that there are only two independent pion--pion scattering
amplitudes, in the limit in which isospin is a good symmetry.

% BANKS-SOURCE: pdf=101 print=91 kind=problem id=7.5
\BanksProblem{7.5}{}
Represent the nucleon at low energy by an isospin doublet field $N$. This field
is invariant under chiral transformations, but since it transforms under the
unbroken isospin subgroup, we must write its Lagrangian with a covariant
derivative. Use $D=4$ Clifford conventions
$\{\gamma^\mu,\gamma^\nu\}=-2\eta^{\mu\nu}$,
$\gamma_5=\ii\gamma^0\gamma^1\gamma^2\gamma^3$, and
$\slashed p=\gamma^\mu p_\mu$:
\begin{equation*}
  \mathcal L=\bar N[-\ii\gamma^\mu(\partial_\mu-A_\mu^a t^a)-m]N.
\end{equation*}
The Goldstone boson fields belong to the coset space
$SU(2)\times SU(2)/SU_V(2)$, where the vector or diagonal subgroup is the
simultaneous action of the same $SU(2)$ transformation in both factors. In
other words we have $(g_L(x),g_R(x))$, with the gauge identification
$(g_L(x),g_R(x))\sim(g_L(x)h(x),g_R(x)h(x))$. Here $g_L$, $g_R$, and $h$ are
$SU(2)$ matrices. Show that the two fields
\begin{equation*}
  A_{\mu L,R}\equiv g_{L,R}^{-1}\partial_\mu g_{L,R}
\end{equation*}
are invariant under the global symmetry and transform as gauge potentials for
the gauge identification. Either one could act as the gauge potential in the
covariant derivative. Use this, and the fact that strong interactions are
invariant under parity (which tells you which combination of $A_{\mu L,R}$
should be substituted for $A_\mu^a$ in the nucleon covariant derivative), to
write down the most general pion--nucleon Lagrangian invariant under the
symmetries of QCD with two massless flavors, and containing at most one
derivative of the pion field in nucleon interactions and two in pion
interactions. In addition to the nucleon mass and pion decay constant, you will
have to introduce one dimensionless coupling $g_A$, determined by the
expectation value of the axial $SU(2)$ current in the nucleon state. This
constant appears in a parity-invariant coupling of the other linear combination
of $A_{\mu L,R}$ to an axial current built from nucleon fields. Use Noether's
theorem to derive the $SU(2)_L\times SU(2)_R$ currents, recalling that in this
formalism only the NGB fields $g_{L,R}$ transform under these symmetries. Now
fix the gauge by carrying out a transformation with $h=g_R^{-1}$. The NGB
field in this gauge is $\Sigma=g_Lg_R^{-1}$ and transforms as
$\Sigma\to V_L^\dagger\Sigma V_R$. Write the Lagrangian in terms of $\Sigma$
(you have to be careful to transform the nucleon field by $h=g_R^{-1}$ to get
this right). Show

% BANKS-SOURCE: pdf=102 print=92 kind=problem id=7.5-cont
that the long-range single-pion exchange force between nuclei is determined by
$g_A$ and $f_\pi$, while pion--pion scattering depends only on $f_\pi$. Now
assume that the $SU(2)_L$ current couples to $W$ bosons. Compute the amplitudes
for neutron decay and charged-pion decay in terms of the strong-interaction
parameters $f_\pi$ and $g_A$, and the Fermi constant $G_F$.

% BANKS-SOURCE: pdf=102 print=92 kind=problem id=7.6
\BanksProblem{7.6}{}
Write the Ward identities for the generating functional of non-abelian
currents, the functional average of
\begin{equation*}
  \left\langle e^{\ii\int\dd^D x\,J_a^\mu(x)A_\mu^a(x)}\right\rangle
  \equiv e^{\ii W(A)}
\end{equation*}
for some global symmetry group $G$ with generators $T^a$. Show that they are
equivalent to the requirement that $W(A)$ be invariant under non-abelian gauge
transformations of $A$. Define the non-standard Legendre transform
\begin{equation*}
  \Gamma(a_\mu^a)+M^2(A_\mu^a-a_\mu^a)(A^{a\mu}-a^{a\mu})=W(A_\mu^a).
\end{equation*}
Describe how the expansion coefficients of $\Gamma$ are related to connected
Green functions. Show that the Ward identities imply that $\Gamma$ is a
gauge-invariant functional of $a_\mu^a$. At low energies this means that
$\Gamma$ is just the usual Yang--Mills action of Chapter 8. Now argue that, if
$G$ is broken to $H$, then we should couple the $a_\mu^a$ field to the NGBs
living in the $G/H$ coset, in a way that respects gauge invariance. Show that
in the Yang--Mills approximation we get a theory of massive vector mesons
interacting with massless NGBs. The vector mass matrix is not just given by
$M^2$. There is also a contribution from the interaction with the NGBs.
Compute it. Apply this formalism to the case of strong interaction chiral
symmetry $G=SU(2)\times SU(2)$. It gives an approximate theory of $\rho$,
$\omega$, and $A_1$ mesons interacting with pions. The theory cannot be
justified in the way that we justify the chiral Lagrangian. Even if we make the
chiral symmetry exact, by setting up and down quark masses to zero, the vector
mesons remain massive, with a mass not much smaller than $4\pi f_\pi$, so there
is no limit in which we can exactly replace $\Gamma$ by the Yang--Mills action.
Nonetheless, this vector-dominance model of the hadronic currents is a useful
approach to problems in strong-interaction physics.

\end{banksproblems}
